\documentclass[11pt]{article}
\usepackage[utf8]{inputenc}
\usepackage[T1]{fontenc}
\usepackage{amsmath,amssymb}
\usepackage{booktabs}
\usepackage{longtable}
\usepackage[margin=1in]{geometry}
\usepackage{hyperref}
\hypersetup{colorlinks=true, linkcolor=blue, urlcolor=blue}

\title{Quantifying the defects in the legacy PUV processing code}
\author{Holden Leslie-Bole\\Scripps Institution of Oceanography, UC San Diego}
\date{27 July 2026}

\begin{document}
\maketitle

\begin{abstract}
\noindent
The legacy PUV processing code (\texttt{beachpeeps/PUV\_Processing}) contains
five identified faults. This note measures what each one does to the derived wave
parameters, so that results produced with that code can be judged individually
rather than dismissed or trusted wholesale. Each legacy choice was implemented as
an isolated switch applied to identical input, making every difference
attributable to one fault. Across 3{,}713 hourly segments from ten instrument
records spanning 5.5--11.7\,m depth, five sites, and $H_s$ to 4.7\,m: significant
wave height and mean wave direction are usable, with stateable offsets; peak
period differs by definition rather than by error and is the largest single
discrepancy between the pipelines; and the directional second moments, from which
radiation stress and alongshore forcing are built, are the products that should
not be relied on. Two results contradicted the expectations formed from reading
the code, and one diagnostic in an earlier version of this comparison overstated
a fault by a factor of three.
\end{abstract}

\section{Why isolate rather than compare pipelines}

Running both pipelines end to end and differencing the outputs confounds all five
faults together with several deliberate design differences. The resulting scatter
shows \emph{that} the pipelines disagree but cannot attribute the disagreement,
and the attribution is the practical question: a result built on $H_s$ may be
sound while one built on radiation stress is not.

Each legacy behaviour was therefore reimplemented from the legacy source as a
switch applied to the same input segments, with everything else held at the
current pipeline's settings. The faults follow the numbering in
\texttt{docs/pipeline\_comparison\_legacy.md}.

\paragraph{D1: auto- and cross-spectra estimated differently.}
Legacy auto-spectra (\texttt{calculate\_fft2}) use a Hanning window with
50\,\% overlap ($\texttt{num} = \lfloor 2n/\texttt{nfft}\rfloor - 1$ ensembles).
Legacy cross-spectra (\texttt{cospec}) use a Hamming window with \emph{zero}
overlap and a hardcoded $\texttt{nfft} = 2400$ carrying the source comment
\emph{``changed from 7200. Why?''}. The two therefore differ in spectral window
\emph{and} in degrees of freedom, while
\begin{equation}
    a_1 = \frac{\mathrm{Re}\,S_{pu}}{\sqrt{S_{pp}\,(S_{uu}+S_{vv})}}
\end{equation}
divides one by the other. Every directional quantity inherits this.

\paragraph{D3: NaN replaced by zero before the FFT.}
\texttt{calculate\_fft2.m:11} and \texttt{cospec.m:22}. The expectation recorded
previously was an $H_s$ deficit proportional to gap fraction.

\paragraph{D4: pressure correction.}
Legacy caps $K_p^2$ at 10; the current pipeline zeros frequencies where
$K_p < 0.1$.

\paragraph{Design differences.} The swell band runs to 0.20\,Hz rather than
0.25\,Hz, and peak period is an energy-weighted centroid rather than the argmax
of $S_\eta$. Neither is a fault.

\paragraph{Not quantified.} D2 (frequency-grid mismatch between \texttt{fm} and
\texttt{f\_co}) and D5 (\texttt{Suv} indexed at line 49, computed at line 107)
are structural faults with no correct magnitude to measure. D5 either errors or
silently reuses a stale workspace variable, so its effect depends on execution
history rather than on the data.

\section{Data}

Ten instrument-records, 3{,}713 hourly segments. Two records (TOR19W and TOR20W,
MOP582, 10--11\,m, 2019--2021) match the station, depth and period of published
work using this code, so the deltas there transfer directly. Five records were
added specifically to reach conditions the first pass could not: $H_s$ to 4.7\,m
(TOR15B), and pressure-gap fractions of 11\,\%, 35\,\% and 85\,\% (TOR14B,
TOR23W/MOP586\_10m, RUBY22/MOP579\_6m).

\section{D1: the estimator mismatch}

\begin{table}[htbp]
  \centering
  \caption{Legacy minus current, directional quantities. $a_1$ and $b_2$ are
  evaluated at the spectral peak and are bounded on $[-1,1]$.}
  \begin{tabular}{@{}lrrr@{}}
    \toprule
    Quantity & Median & IQR & 95th pct.\ of $|\Delta|$ \\
    \midrule
    Mean direction $D_p$ (deg)      & $-0.07$ & $-1.40$ to $+1.40$ & 6.6 \\
    Spread, peak bin (deg)          & $-4.32$ & $-7.78$ to $+29.03$ & 50.3 \\
    Spread, band-averaged (deg)     & $+0.26$ & $-9.45$ to $+6.24$ & \textbf{17.1} \\
    $a_1$ at peak                   & $-0.001$ & $-0.043$ to $+0.040$ & 0.177 \\
    $b_2$ at peak                   & $+0.002$ & $-0.080$ to $+0.085$ & \textbf{0.326} \\
    \bottomrule
  \end{tabular}
\end{table}

Mean direction is essentially unbiased in the median because the fault injects
variance rather than a systematic offset, though individual hours reach
6.6$^\circ$. The second moment does not fare as well: $b_2$ has a
95th-percentile error of 0.326 on a quantity that spans only $[-1,1]$. Radiation
stress
\begin{equation}
    S_{xy} = \rho g \int S(f)\, \frac{c_g}{c}\, \frac{b_2(f)}{2}\, df
\end{equation}
is built directly from $b_2$ and inherits that error in full, which makes legacy
alongshore forcing the least reliable product of that pipeline.

\subsection*{A correction to the earlier version of this comparison}

The peak-bin spread figure above (95th percentile 50.3$^\circ$) overstates the
fault by roughly a factor of three, and an earlier draft of this comparison
reported it without qualification. Spread computed at a single bin from
\begin{equation}
    \sigma_1 = \sqrt{2\left(1 - r_1\right)}, \qquad r_1 = \sqrt{a_1^2 + b_1^2}
\end{equation}
is dominated by the degrees-of-freedom difference between the two estimators
rather than by the window mismatch: with zero-overlap cross-spectra the peak-bin
$r_1$ is noisy and occasionally exceeds unity, where it must be clamped.
Averaging the complex first moment over the swell band before taking its
modulus,
\begin{equation}
    \bar r_1 = \left| \frac{\sum_f S(f)\,[a_1(f) + i\,b_1(f)]}{\sum_f S(f)} \right|,
\end{equation}
suppresses that noise and isolates the window mismatch. The band-averaged
statistic is the defensible one and is closer to the form in general use
following Kuik et al.\ (1988). The fault is real, and 17$^\circ$ at the 95th
percentile is still large, but it is not 50$^\circ$.

\subsection*{Dependence on conditions}

\begin{table}[htbp]
  \centering
  \caption{Median $|\Delta|$ by significant wave height.}
  \begin{tabular}{@{}lrrrr@{}}
    \toprule
    $H_s$ (m) & $n$ & $|\Delta D_p|$ (deg) & $|\Delta\sigma_1|$ (deg) & $|\Delta b_2|$ \\
    \midrule
    $<0.5$    &  374 & 1.71 & 6.22 & 0.129 \\
    0.5--1.0  & 1816 & 1.48 & 7.07 & 0.095 \\
    1.0--1.5  &  889 & 1.34 & 8.14 & 0.066 \\
    1.5--2.0  &  359 & 1.33 & 8.63 & 0.059 \\
    2.0--3.0  &  247 & 1.11 & 7.79 & 0.065 \\
    $>3.0$    &   28 & 0.83 & 6.62 & 0.099 \\
    \bottomrule
  \end{tabular}
\end{table}

The direction error falls monotonically with wave height, from 1.71$^\circ$ below
$H_s = 0.5$\,m to 0.83$^\circ$ above 3\,m, and $|\Delta b_2|$ falls by half over
the same range. This is the signature expected of a variance-driven fault: at
higher signal-to-noise the two estimators converge. \textbf{Large waves are the
regime in which the legacy directional output is most trustworthy, not least.}
The apparent uptick in $|\Delta b_2|$ in the highest bin rests on 28 segments and
should not be read as a reversal.

The error grows with directional spread instead, from a median
$|\Delta D_p|$ of 1.35$^\circ$ below 15$^\circ$ spread to 2.79$^\circ$ at
20--25$^\circ$: broader seas have lower coherence per frequency, so the
estimator difference matters more.

\section{D3: the NaN handling, and why it does not bite}

This fault could not be made to act as documented, and the reason is more useful
than the number would have been.

Records were selected specifically for pressure gaps --- 11\,\%, 35\,\% and
85\,\% NaN in the raw pressure. After the standard validity gate, essentially
every surviving segment has a pressure-gap fraction below 1\,\%. The segments in
which zeroing would matter do not reach the spectral estimator, because the
existing quality control rejects them first.

Where small gaps do survive, the fractional change in $H_s$ is erratic rather
than a deficit: the median is zero, the interquartile range reaches $+0.19$, and
gap fraction does not predict the change ($\rho = -0.06$, $n = 1152$). The
predicted mechanism --- variance removed in proportion to the gap --- competes
with a second effect in the opposite direction, since zeroing a segment of an
oscillating record introduces step discontinuities that inject broadband energy.

The gaps in these records are overwhelmingly in \emph{velocity}, not pressure:
8.8\,\% NaN in $U$ and $V$ on TOR19W against $2\times10^{-5}$ in pressure, and
24\,\% on COR16B. So the fault does not act on $H_s$ at all. It zeros roughly one
velocity sample in eleven before the FFT, which removes velocity variance and
therefore acts on the directional coefficients and on $U_b$ --- compounding D1 on
exactly the quantities D1 already damages. The original write-up attributed this
fault to the wrong variable.

\textbf{The protection here is the validity gate, not the NaN handling.} Any
reprocessing that relaxes that gate would expose this fault.

\section{D4: the pressure-correction rule}

\begin{table}[htbp]
  \centering
  \caption{Median $|\Delta H_s|/H_s$ from capping $K_p^2$ at 10 rather than
  zeroing frequencies below $K_p = 0.1$.}
  \begin{tabular}{@{}lrr@{}}
    \toprule
    Depth (m) & $n$ & Median $|\Delta H_s|/H_s$ \\
    \midrule
    $<6$    &  339 & 0.000 \\
    6--8    &  867 & 0.000 \\
    8--10   & 1387 & 0.005 \\
    $>10$   & 1120 & \textbf{0.015} \\
    \bottomrule
  \end{tabular}
\end{table}

This runs opposite to the expectation recorded from the code reading, which
argued from noise amplification being permitted in shallow water. Measured on
$H_s$, the two rules diverge in \emph{deeper} water: there $K_p$ falls below the
cutoff across a wider portion of the sea-swell band, so the capped and zeroed
treatments disagree over more frequencies. The effect is 1.5\,\% above 10\,m and
negligible below 8\,m.

The shallow-water concern in the original reading is not thereby dismissed. It
was about noise admitted at high frequency, which affects the spectral tail and
any quantity weighted toward it; $H_s$ is an integral dominated by the peak and
is a poor detector of it.

\section{Design differences, which are larger than most of the faults}

\paragraph{Swell-band upper limit.} Truncating at 0.20\,Hz rather than 0.25\,Hz
makes legacy $H_s^{SS}$ 2.9\,\% low in the median, rising to 7.8\,\% for
$T_p < 8$\,s and falling to 2.5\,\% for $T_p > 16$\,s. A wind-sea effect,
negligible in swell.

\paragraph{Peak period.} The energy-weighted centroid sits 2.73\,s below the
argmax definition in the median, with an interquartile range of $-4.29$ to
$-1.47$\,s and a 95th percentile of 7.11\,s. \textbf{This is the largest single
difference between the pipelines and it is not an error in either.} The two
quantities are different statistics of the same spectrum. A legacy $T_p$
compared against a model $T_p$, or against a value from another study, is a
comparison of definitions unless the centroid is stated explicitly.

\section{Summary}

\begin{longtable}{@{}p{0.24\textwidth}p{0.30\textwidth}p{0.38\textwidth}@{}}
  \toprule
  Quantity & Status & Qualification \\
  \midrule
  \endhead
  $H_s$ & Usable & 2.9\,\% low from the band limit; a further 1.5\,\% above
  10\,m depth from the $K_p$ rule. Both are offsets, not errors. \\
  \addlinespace
  Mean direction $D_p$ & Usable & Median bias below 0.1$^\circ$; individual hours
  reach 6.6$^\circ$. Best at large $H_s$, worst in broad seas. \\
  \addlinespace
  $T_p$ & Different, not wrong & 2.73\,s below an argmax $T_p$. State the
  definition before comparing with anything. \\
  \addlinespace
  Directional spread & Distrust & 17$^\circ$ at the 95th percentile,
  band-averaged. \\
  \addlinespace
  $b_2$, $S_{xy}$, alongshore forcing & Do not use & 95th-percentile $b_2$ error
  of 0.33 on a range of 2, with D3 compounding it through the velocity channel. \\
  \bottomrule
\end{longtable}

\section{Limits of this comparison}

Each fault was emulated from the legacy source rather than by executing that
pipeline end to end, so interactions between faults would not appear here. Ten
records at five sites is a narrow span: only 28 segments exceed $H_s = 3$\,m, and
no valid segment carries a pressure-gap fraction above 1\,\%, so the two regimes
this note set out to probe are covered thinly at one end and not at all at the
other. D2 and D5 remain unquantified.

The band-averaged spread reported here should be preferred over the peak-bin
figure, including over the peak-bin figure that appeared in the earlier version
of this comparison.

\vspace{1em}
\noindent\rule{\textwidth}{0.4pt}\\
\small
Reproduce with \texttt{validation/legacy\_defect\_isolation.m}; results in
\texttt{outputs/validation/legacy\_defect\_isolation.mat}. Legacy source:
\texttt{Beach\_Change\_Observation/Vector/PUVs/PUV\_Processing-main/Level2\_QC/}.
Prepared as internal documentation of why the legacy codes are deprecated.

\end{document}
